\documentclass[11pt,a4paper]{article}
\usepackage[margin=23mm]{geometry}
\usepackage{amsmath,amssymb,amsthm,hyperref}
\newtheorem{lemma}{Lemma}
\newcommand{\gen}[1]{\langle#1\rangle}
\title{Small endpoints in the single-swap block case}
\author{Human Compression project: further bounded L044 proof}
\date{}
\begin{document}\maketitle
We retain the hypotheses and notation of the accepted note
\emph{Exact block and orbit deficits at a single crossing}.
In particular $G\leq A_n$ is a transitive string C-group, every
maximal omission is intransitive, and the unique eligible edge
of colour $i$ joins the two $G_i$-orbits. Each orbit now has
two blocks, of sizes $k_1,k_2\geq4$, and $\alpha=1$.
Put $s=\rho_i$. No assertion below uses a projected intersection
property before the projection is proved faithful.

\begin{lemma}[A four-row endpoint forces redundancy]\label{four}
In the stated setting,
\[
 r\leq\left\lfloor\frac{n-1}{2}\right\rfloor.
\]
Thus the accepted single-swap conclusion extends from
$k_1,k_2\geq5$ to $k_1,k_2\geq4$.
\end{lemma}
\begin{proof}
Suppose the bound fails. The exact deficit identity gives
$\beta=\gamma=0$, $r=k_1+k_2$, and a single simultaneous
fixed-point-free block swap $\tau=\rho_{i-1}$, after reversing
the order if necessary. The proof of Lemma~14 of the accepted
note, up to its endpoint analysis, applies for $k_1,k_2\geq4$:
the faithful normalized group has the form
\[
 H=S_a\times S_{b-1},\qquad \{a,b\}=\{k_1,k_2\}.
\]
Here its character is trivial, its right string interval is
the diagonal natural $S_a$ moving its row anchor, and its left
interval is the diagonal $S_{b-1}$ fixing its row anchor.
Both strings have maximal independent length. All symmetric
factors used to separate these intervals are centerless,
since $a,b\geq4$; this is the only lower-degree restriction in
that part of the argument. The generator $x=\rho_{i-2}$ exists,
fixes the two crossing endpoints, preserves all four blocks,
centralizes the right interval and the left endpoint omission,
and must connect the fixed row of the left interval to its
other rows. The accepted argument excludes $b\geq5$ without
requiring $a\geq5$: the right natural $S_a$ has trivial
centralizer for every $a\geq3$.

It remains to take $b=4$. Write the rows of its two blocks as
$0,1,2,3$ and $0',1',2',3'$, with $0$ the crossing endpoint.
Its left interval has two generators, necessarily two distinct
transpositions on the three nonanchor rows. Relabel them as
\[
 p=(1\ 2)(1'\ 2'),\qquad
 q=(2\ 3)(2'\ 3').
\]
They occur in the order $p,q,x,\tau,s$. As $x$ centralizes
the complete right-hand $S_a$, it is the identity on that
original orbit. On the present orbit it centralizes $p$.
Its first-block action fixes $0$, so is $(1\ 2)^\epsilon$.
Its second-block action is
$(1'\ 2')^\delta(0'\ 3')^\eta$.
If $\eta=0$, the pair $\{0,0'\}$ is invariant under
$p,q,x,\tau$ and every remaining generator of $G_i$,
contradicting this orbit's transitivity. Thus $\eta=1$.
Evenness of $x$ gives $\epsilon+\delta=1$ modulo two.
There are exactly two possibilities:
\[
 x=(1\ 2)(0'\ 3')\quad\hbox{or}\quad
 x=(1'\ 2')(0'\ 3').                                      \tag{1}
\]

The subgroup $\gen{p,q,x}$ has, on these eight points, orbits
$\{0\}$, $\{1,2,3\}$ and $\{0',1',2',3'\}$.
Its projections on the latter two are the natural $S_3$ and
$S_4$: the second contains $(1'\ 2'),(2'\ 3')$ and,
after multiplication by $(1'\ 2')$ if necessary, $(0'\ 3')$.
These two natural actions have trivial centralizer and
different orbit sizes. Since $s$ commutes with $p,q,x$, it
fixes those seven points individually.

Write the rows of the other orbit as
$A_0,\ldots,A_{a-1}$ and $A'_0,\ldots,A'_{a-1}$, with
$A_0$ the other crossing endpoint. Only the first generator
of the maximal right-hand $S_a$ string moves its row anchor.
The anchored classification in Lemma~6 of the accepted note
therefore shows that its tail is the natural diagonal
$S_{a-1}$ on rows $1,\ldots,a-1$. This includes $a=4$:
the first generator may have its extra disjoint transposition,
but the last two still generate the full point stabilizer
$S_3$. The outer degree-six action is excluded by the fixed
anchor, exactly as in the accepted lemma.

The generator $s$ centralizes this tail. Its only possible
action on the two $(a-1)$-point orbits is to interchange them
row by row, or to fix both pointwise. On the remaining three
points $A_0,A'_0,0$, the unique crossing gives $(A_0\ 0)$,
so $A'_0$ is fixed. Fixing the other rows would make $s$ one
odd transposition. Hence
\[
 s=(A_0\ 0)\prod_{j=1}^{a-1}(A_j\ A'_j),\qquad
 \tau=\prod_{j=0}^{a-1}(A_j\ A'_j)
            (0\ 0')(1\ 1')(2\ 2')(3\ 3').                 \tag{2}
\]
In particular $a$ is even, consistently with the sign of
the fixed-point-free $\tau$.

With right actions and products applied from left to right,
both alternatives in (1) satisfy the exact identity
\[
                 p=s(\tau s x)^5\tau x.                    \tag{3}
\]
For completeness this identity can be checked without a
group-order calculation. On the ten points
$A_0,A'_0,0,1,2,3,0',1',2',3'$, the permutation $\tau s x$
has cycles
\[
 (A_0\ A'_0\ 0\ 3'\ 3\ 0')
 \quad\hbox{and}\quad
 \begin{cases}
 (1\ 1'\ 2\ 2'),&x=(1\ 2)(0'\ 3'),\\
 (1\ 2'\ 2\ 1'),&x=(1'\ 2')(0'\ 3').
 \end{cases}
\]
Its fifth power reverses the six-cycle and leaves the
four-cycle as displayed. Multiplication as in (3) gives
$(1\ 2)(1'\ 2')$ in each case. On every remaining pair
$\{A_j,A'_j\}$, both $s$ and $\tau$ exchange the two points
and $x$ is the identity. The word in (3) is therefore the
identity there. This proves (3) on all points, for every
permitted $a$.

Equation (3) makes the distinguished generator $p$ redundant
in the full marking. This contradicts IP, or just independence,
and completes the proof.
\end{proof}

\begin{lemma}[A three-row endpoint]\label{three}
The conclusion also holds with the initial restriction
$k_1,k_2\geq4$ relaxed to $k_1,k_2\geq3$.
\end{lemma}
\begin{proof}
Only a block size three remains after Lemma~\ref{four}.
The faithful normalization from the preceding proof still
gives $H=S_a\times S_{b-1}$ with a trivial character, where
$a$ is the row orbit whose anchor is moved by the right-hand
generator $c$. Indeed the point stabilizer $S_{a-1}$ normally
generates $S_a$ already for $a=3$, and the two-anchor
independent-length argument also applies at this endpoint.
If $b\geq4$, the two factors remain centerless and the
accepted endpoint analysis applies as before. The case $b=4$
with $a=3$ is excluded immediately because $\tau$ would
be the product of seven transpositions. Thus assume $b=3$.
Evenness of $\tau$ forces $a$ odd.

The left interval must be nonempty: otherwise $x$ centralizes
all of $H$, fixes the unique fixed row of the natural $S_2$
on each three-point block, and leaves its two-point pair
invariant under $G_i$. The right interval contains $c$.
The projections of the two commuting interval factors onto
$S_a$ are commuting normal subgroups whose product is $S_a$.
Their intersection is central, hence trivial. Since $S_a$
is directly indecomposable, one projection is trivial; it
must be the left one, since the right interval moves the
anchor. The left interval is consequently the full central
$C_2$ of $H=S_a\times C_2$, generated by one element $p$.
The right interval is a complement and hence the graph of
the trivial or sign homomorphism $S_a\to C_2$.

The sign graph is impossible. The generator $x$ centralizes
the right interval, so is the identity on the original orbit
with row group $S_a$. On the three-row orbit it then
centralizes the nonidentity $C_2$ action occurring in an odd
right-interval element. This fixes its unique fixed row in
each block, leaving the anchor pair invariant under $G_i$.
Thus the right interval acts trivially on this orbit.

Label its blocks $0,1,2$ and $0',1',2'$. Then
$p=(1\ 2)(1'\ 2')$, and $x$ fixes the whole other original
orbit. On the first three-point block it fixes $0$, and on
the second it must move $0'$ to make the orbit transitive.
Evenness and a simultaneous exchange of names $1,2$ give
\[
 x=(1\ 2)(0'\ 1').                                      \tag{4}
\]
The subgroup $\gen{p,x}$ has a two-point orbit $\{1,2\}$
and a three-point orbit $\{0',1',2'\}$, on which its images
are $S_2$ and $S_3$. Thus $s$, which commutes with both
generators, fixes the latter orbit pointwise and on the former
may perform $(1\ 2)$.

First take $a\geq5$. The anchored maximal right interval is
the ordinary diagonal Coxeter string, so put
$c=(A_0\ A_1)(A'_0\ A'_1)$ for its first generator.
Its tail is diagonal $S_{a-1}$ on all other rows.
Consequently, for some $\epsilon\in\{0,1\}$,
\[
 \tau=\prod_{j=0}^{a-1}(A_j\ A'_j)(0\ 0')(1\ 1')(2\ 2'),
 \quad
 s=(A_0\ 0)(1\ 2)
       \left(\prod_{j=1}^{a-1}(A_j\ A'_j)\right)^\epsilon. \tag{5}
\]
Here $(1\ 2)$ is forced by parity: $a-1$ is even, so either
allowed action on those remaining rows is even, while the
crossing transposition is odd. Both choices satisfy
\[
                         p=(c s\tau s)^3\tau.             \tag{6}
\]
To verify it, on the ten points over $A_0,A_1$ and the
three-row orbit, $c s\tau s$ has cycles
\[
 (A_0\ A'_1\ 0\ A'_0\ A_1\ 0'),\qquad
 (1\ 2'),\qquad(2\ 1').
\]
Cubing and multiplying by $\tau$ gives $p$. On every
remaining row pair, $c s\tau s$ is $\tau$, so cubing and
multiplying by $\tau$ is the identity. This proves (6)
uniformly and contradicts independence of $p$.

For clarity the degree-twelve endpoint $a=b=3$ requires
one additional centralizer check. The last generator is
$t=(A_1\ A_2)(A'_1\ A'_2)$, and its centralizer on these
four points is $D_8$. Write $s=(A_0\ 0)h(1\ 2)^\eta$.
Its evenness and involutivity leave precisely
\[
\begin{array}{c|l}
\eta&h\\\hline
1&1,\quad (A_1\ A_2)(A'_1\ A'_2),\quad
 (A_1\ A'_1)(A_2\ A'_2),\quad
 (A_1\ A'_2)(A_2\ A'_1),\\
0&(A_1\ A_2),\quad (A'_1\ A'_2).
\end{array}                                               \tag{7}
\]
For the first four choices, $h$ commutes with the restriction
of $\tau$ and fixes the crossing points. Hence $s\tau s$
and the cycle verification of (6) are unchanged.
For either of the last two choices, direct multiplication
using (4), (7) and the displayed $\tau$ gives instead
\[
                p=\tau s(\tau x\tau s x)^3\tau x.          \tag{8}
\]
Indeed $\tau x\tau s x$ has disjoint cycles
$(A_0\ 0\ 2\ 1)$, $(0'\ 1'\ 2')$, and $h$.
Its cube reverses the four-cycle, kills the three-cycle and
retains $h$, giving (8) by multiplication. Thus this is a
literal identity on the twelve explicitly named points.
All six possibilities again make $p$ redundant.
This completes the endpoint and the proof.
\end{proof}

\paragraph{Scope and attribution.}
The source branch is Section~5.1 of Cameron--Fernandes--Leemans--Mixer,
\emph{Highest rank of a polytope for $A_n$}, arXiv:1605.09173v1.
The simultaneous-swap and row-normalization framework is the
one proved in the accepted block-deficit note, with its full
L047/Whiston/extremal dependencies retained. The explicit
uniform redundancy words (3), (6) and (8) are the new arguments here.
A bounded exploratory permutation calculation suggested that
word; the displayed cycle verification is its complete proof,
not a sample-based extrapolation or an exhaustive-census premise.
The cases with a block size two, the two-independent-swap
case, and the remaining L044 principal obligations are not
settled by this note. There is no new original, conjecture or
formalization claim.
\end{document}
